\chapter*{前言}
\addcontentsline{toc}{chapter}{前言}
\sourcepage{xi--xii}
代数拓扑是现代数学的基石之一。每一位从事数学工作的人都应当至少对这门学科有所了解。James R. Munkres 的《代数拓扑基础》是一部提供这种入门介绍的经典著作。如今，这本书问世已有四十年，需要加以更新。本次新版正是为了实现这一目的。

当然，我们保留了原书的基本特色。大部分记号和术语保持不变，第一版的所有核心主题也都予以保留。不过，新版作出了一些有益的增补和修改。

与第一版不同，我们在讨论基本群的一节中介绍了同伦论。无环模型一节经过全面重写，以便读者更容易理解。书末集中收录了正文中提到的各位人物的一些简短生平资料。此外，还列出了一份简短的补充阅读书目，其中大多数书籍出版于第一版之后。另一项新增内容是十分详尽的符号索引；全书的主题索引现在也更加详尽。最后，我们增加了相当数量的习题，其中多数属于巩固练习。

\section*{本书的组织结构}
前两章以单纯同调为基础；在各种同调理论中，它最为具体，因此学生应当比较容易掌握。第三章介绍更抽象的内容，包括 Eilenberg--Steenrod 公理。第四章介绍奇异同调论及其许多有趣的应用，并把 CW 复形作为计算工具引入。对于一学期的课程，前四章是合适的内容。

标有星号的各节，在这一学期课程中可以作为选学内容。当然，在作者看来，那些应用十分有趣，因此即使属于选学内容，教师仍可能选择讲授其中一部分或全部。

第五章介绍上同调论，并最终讨论上同调环。第六章引入张量积，并进一步讨论第一章某个选学小节中提到的任意系数同调。虽然前六章的内容很可能超出一学期课程的容量，但它们构成了扎实的核心内容，应当适合相当广泛的数学工作者，包括学生和资深数学家。

对于两学期的课程，整本书都可以作为教学内容。当然，教师在设计合适的课程时需要有所取舍。对一学期课程而言属于选学的各节，例如“范畴与函子”，对于两学期课程可能是必要的。最后两章彼此独立，因此，如果第二学期接近结束时必须删减一些内容，也可以灵活处理。

具体来说，第七章讨论同调代数，包括各种情形下的泛系数定理和 Künneth 定理。第八章讨论流形，特别是与 Poincaré、Lefschetz、Alexander 和 Pontryagin 的名字相联系的各种对偶定理。为研究 Pontryagin 对偶，本章还引入 Čech 上同调。

如果教师希望减少讲授单纯同调所花的时间，可以略去第二章。这样一来，单纯同调群的拓扑不变性就不再直接证明，而必须通过奇异同调群的不变性，以及单纯同调与奇异同调之间的同构来获得。

\section*{预备知识}
我们假定学生对一般拓扑已有一定了解，包括紧致性、连通性和商空间。书中将用到正规空间上的 Tietze 扩张定理，而商空间尤其重要。我们将在适当之处复习商空间。

我们也希望学生具有一定的代数知识，特别是熟悉群、商群和同态，以及环、域和向量空间。书中不需要用到特别深入的定理。与商空间一样，我们将在适当之处复习基本的代数结果。

\chapter*{致谢}
\addcontentsline{toc}{chapter}{致谢}
\sourcepage{xiii}
我们感谢加利福尼亚大学伯克利分校的 Carlo H. Séquin 提供封面插图。该图展示了两个 Möbius 带如何拼合成一个 Klein 瓶。

感谢编辑 Robert Ross 给予我们的支持与鼓励。我们期待收到读者的反馈。

\begin{flushright}
Steven G. Krantz\\Harold R. Parks
\end{flushright}
